\documentclass[11pt,a4paper]{article}
\usepackage[utf8]{inputenc}
\usepackage{amsmath}
\usepackage{amsfonts}
\usepackage{amssymb}
\usepackage{geometry}

\geometry{top=2.5cm, bottom=2.5cm, left=2.5cm, right=2.5cm}

\title{\textbf{Penyelesaian Luas Permukaan Benda Putar Kurva Parametrik}}
\author{Ahmad Fakhrul Bawani}
\date{}

\begin{document}

\maketitle

\section{Deskripsi Masalah}
Diberikan sebuah kurva parametrik yang didefinisikan oleh persamaan:
\begin{align}
  x &= \frac{1}{2}\cos(2t) \\
  y &= 3 + \frac{1}{2}\sin(2t)
\end{align}
pada rentang parameter $0 \le t \le \frac{\pi}{2}$. Hitunglah luas permukaan benda putar (\textit{surface area}) yang dihasilkan jika kurva tersebut diputar sejauh $360^\circ$ mengelilingi \textbf{sumbu-}$x$.

\section{Formulasi Rumus Luas Permukaan Benda Putar}
Luas permukaan benda putar yang dihasilkan dari rotasi suatu kurva parametrik mengelilingi sumbu-$x$ dapat dicari dengan rumus integral:
\begin{equation}
  S = \int_{\alpha}^{\beta} 2\pi y \sqrt{\left(\frac{dx}{dt}\right)^2 + \left(\frac{dy}{dt}\right)^2} \, dt
\end{equation}

\section{Langkah-Langkah Perhitungan}

\subsection{Menentukan Turunan Komponen Parametrik}
Pertama, cari turunan pertama dari masing-masing komponen variabel terhadap parameter $t$:
\begin{align*}
  \frac{dx}{dt} &= \frac{1}{2} \cdot \Big(-\sin(2t) \cdot 2\Big) = -\sin(2t) \\
  \frac{dy}{dt} &= 0 + \frac{1}{2} \cdot \Big(\cos(2t) \cdot 2\Big) = \cos(2t)
\end{align*}

\subsection{Menyederhanakan Panjang Diferensial Busur}
Lakukan penjumlahan kuadrat turunan untuk menyederhanakan ekspresi di bawah tanda akar dengan menggunakan identitas dasar Pythagoras trigonometri $\sin^2\theta + \cos^2\theta = 1$:
\begin{align*}
  \left(\frac{dx}{dt}\right)^2 + \left(\frac{dy}{dt}\right)^2 &= \Big(-\sin(2t)\Big)^2 + \Big(\cos(2t)\Big)^2 \\
  &= \sin^2(2t) + \cos^2(2t) = 1
\end{align*}
Maka bagian akarnya bernilai konstan:
\begin{equation*}
  \sqrt{\left(\frac{dx}{dt}\right)^2 + \left(\frac{dy}{dt}\right)^2} = \sqrt{1} = 1
\end{equation*}

\subsection{Evaluasi Integral Luas Permukaan}
Substitusikan fungsi $y$, batas nilai $t$ ($\alpha = 0$ dan $\beta = \frac{\pi}{2}$), serta elemen akar yang telah disederhanakan ke dalam formula utama:
\begin{align*}
  S &= \int_{0}^{\frac{\pi}{2}} 2\pi \left( 3 + \frac{1}{2}\sin(2t) \right) \cdot (1) \, dt \\
  S &= 2\pi \int_{0}^{\frac{\pi}{2}} \left( 3 + \frac{1}{2}\sin(2t) \right) \, dt
\end{align*}

Lakukan proses integrasi komponen demi komponen:
\begin{align*}
  S &= 2\pi \left[ 3t - \frac{1}{4}\cos(2t) \right]_{0}^{\frac{\pi}{2}}
\end{align*}

Substitusikan nilai batas atas dan batas bawah ke dalam fungsi hasil integral tersebut:
\begin{align*}
  S &= 2\pi \left[ \left( 3\left(\frac{\pi}{2}\right) - \frac{1}{4}\cos\left(2 \cdot \frac{\pi}{2}\right) \right) - \left( 3(0) - \frac{1}{4}\cos(0) \right) \right] \\
  S &= 2\pi \left[ \left( \frac{3\pi}{2} - \frac{1}{4}(-1) \right) - \left( 0 - \frac{1}{4}(1) \right) \right] \\
  S &= 2\pi \left[ \frac{3\pi}{2} + \frac{1}{4} + \frac{1}{4} \right] \\
  S &= 2\pi \left[ \frac{3\pi}{2} + \frac{1}{2} \right]
\end{align*}

Kalikan konstanta luar $2\pi$ ke dalam kurung siku untuk mendapatkan bentuk penyelesaian akhir:
\begin{align*}
  S &= 3\pi^2 + \pi = \pi(3\pi + 1)
\end{align*}

\section{Kesimpulan}
Luas permukaan benda putar yang terbentuk dari perputaran kurva parametrik tersebut terhadap sumbu-$x$ adalah tepat \textbf{$\pi(3\pi + 1)$} atau \textbf{$3\pi^2 + \pi$} satuan luas.

\end{document}